% ------------------------------------------------------------------------------
% Chapter 1
% ------------------------------------------------------------------------------
\chapter{Design Overview} 
\label{cha:design_overview}

There are 2 font-end modules under test: Fetch and Decode.

\section{Fetch}

The fetch unit is responsible for supplying instructions to the pipeline.
It maintains a program counter (PC), updates it according to control signals, and outputs the instruction stored at the current PC address.
The unit supports sequential execution andbranch redirection.
Instruction memory is modeled internally as a simple array initialized with dummy values.

The fetch unit maintains an internal program counter (PC) that determines which instruction is supplied to the pipeline.
On every rising edge of the clock, the PC is updated according to a strict priority of control signals.
If reset is asserted, the PC isforced to zero regardless of any other inputs.
If reset is not active and a branch request is signaled, the PC is immediately loaded with the branch target address.
In the absence of a branch, if the sequential enable signal is asserted, the PC advances by 2, assuming 16-bit instructions and byte addressing.
When neither branch nor sequential enable is active, the PC simply retains its previous value.
The next PC value is computed combinationally based on control signals, ensuring that both the PC and instruction outputs are updated synchronously and consistently on each clock edge.
The instruction output reflects the contents of the memory location that will be addressed by the PC in the next cycle.
This ensures that the instruction is fetched ahead of execution and remains consistent with the updated PC value.
Memory is modeled as a 256‑entry array of sixteen‑bit words, each initialized with a dummy value equal to its address index.
This initialization strategy allows for easy verification of PC behavior, as the instruction output directly reflects the PC value.
Branch requests take precedence over sequential increments, so if both signals are asserted in the same cycle, the branch target address is chosen.
The PC must remain within the legal range of zero to 255 at all times, and any attempt to access outside this range is considered invalid behavior.
Together, these rules define the functional behavior of the fetch unit: reset initializes the PC to zero, branch requests redirect execution, sequential enable advances the PC, and the instruction output is always consistent with the memory contents at thecurrent PC.

\subsection{Interface}

\begin{table}[h!]
    \centering
    \begin{tabular}{|l|c|c|l|}
        \hline
        \textbf{Signal}         & \textbf{Direction} & \textbf{Width} & \textbf{Description}            \\ \hline
        \texttt{clk}            & Input              & 1              & Clock signal                    \\ \hline
        \texttt{rst\_n}         & Input              & 1              & Active-low synchronous reset    \\ \hline
        \texttt{pc\_en}         & Input              & 1              & Program counter enable          \\ \hline
        \texttt{branch\_en}     & Input              & 1              & Branching enable                \\ \hline
        \texttt{branch\_addr}   & Input              & 8              & Branching address               \\ \hline
        \texttt{instr}          & Output             & 16             & Instruction fetched             \\ \hline
    \end{tabular}
    \caption{Register File signals}

\end{table}

\section{Decode}

The decode stage is a pipeline block that interprets 16 bit instructions received from the fetch stage.
It extracts the opcode, destination register, source register, and immediate field, and prepares them for the execute stage.
To model realistic pipeline behavior, the decode stage introduces a fixed two cycle latency between instruction acceptance and completion.
It also performs simple hazard detection: if the destination register of the previous instruction matches the source register of the current one, the stage asserts a stall signal to delay acceptance until the hazard clears.
The stage receives instructions from fetch via an instr\_validsignal and signals completion with a single-cycle decode\_donepulse.

\subsection{Interface}

\begin{table}[h!]
    \centering
    \begin{tabular}{|l|c|c|l|}
        \hline
        \textbf{Signal}         & \textbf{Direction} & \textbf{Width} & \textbf{Description}                    \\ \hline
        \texttt{clk}            & Input              & 1              & Clock signal                            \\ \hline
        \texttt{rst\_n}         & Input              & 1              & Active-low synchronous reset            \\ \hline
        \texttt{instr\_valid}   & Input              & 1              & notify succesful fetch                  \\ \hline
        \texttt{instr}          & Input              & 16             & Instruction fetched                     \\ \hline
        \texttt{decode\_done}   & Output             & 1              & Pulse signaling the succesful decode    \\ \hline
        \texttt{opcode}         & Output             & 4              & decoded opcode                          \\ \hline
        \texttt{rd}             & Output             & 4              & decoded destination register            \\ \hline
        \texttt{rs}             & Output             & 4              & decoded source register                 \\ \hline
        \texttt{imm}            & Output             & 4              & decoded immediate value                 \\ \hline
        \texttt{hazard\_stall}  & Output             & 1              & data hazard stalls                      \\ \hline
    \end{tabular}
    \caption{Register File signals}
\end{table}

\subsection{Behavior}

\begin{itemize}
    \item \textbf{Instruction acceptance}
    \begin{itemize}
        \item An instruction is accepted when instr\_valid=1 and hazard\_stall=0 in the nextcycle.
        \item On acceptance, the fields (opcode, rd, rs, imm) must capture the corresponding bits of instr
    \end{itemize}

    \item \textbf{Decode latency}
    \begin{itemize}
        \item After acceptance in cycle N, decode\_donemust assert exactly in cycle N+2.
        \item decode\_done must be a single cycle pulse
    \end{itemize}

    \item \textbf{Hazard detection}
    \begin{itemize}
        \item A hazard occurs when the destination register (rd) of the previously accepted instruction matches the source register (rs) of the current instruction.
        \item When a hazard is detected:
        \begin{itemize}
            \item hazard\_stallmust assert.
            \item The current instruction is not accepted until the stall clears.
            \item decode\_done must not assert during the stall window.
        \end{itemize}
    \end{itemize}

    \item \textbf{Field stability}
    \begin{itemize}
        \item While instr\_valid=1 and hazard\_stall=1, the decoded fields (opcode, rd, rs, imm) must retain their previous values and only update when the instruction is accepted.
    \end{itemize}

    \item \textbf{Back-to-back instructions}
    \begin{itemize}
        \item If instr\_valid is high in consecutive cycles and no hazard occurs, each instruction is accepted immediately.
        \item Each accepted instruction produces a decode\_donepulse two cycles later, possibly resulting in consecutive decode\_donepulses.
    \end{itemize}

    \item \textbf{reset behavior}
    \begin{itemize}
        \item When rst\_n=0:
         \begin{itemize}
            \item All outputsmust be cleared to zero.
            \item No instruction is accepted or completed until reset is released.
        \end{itemize}       
    \end{itemize}
\end{itemize}
